\documentclass[11pt]{article}
\usepackage[margin=1in]{geometry}
\usepackage{amsmath, amssymb, amsthm}
\usepackage{algorithm}
\usepackage{algpseudocode}
\usepackage{booktabs}
\usepackage{hyperref}

\newtheorem{theorem}{Theorem}
\newtheorem{lemma}[theorem]{Lemma}
\newtheorem{proposition}[theorem]{Proposition}

\title{Project Euler Problem 454:\\Diophantine Reciprocals III}
\author{Solution Document}
\date{}

\begin{document}
\maketitle

\section{Problem Statement}

In the equation
\[
\frac{1}{x} + \frac{1}{y} = \frac{1}{n},
\]
where $x$, $y$, and $n$ are positive integers, define $F(L)$ as the number of
solutions satisfying $x < y \leq L$.

Given: $F(15) = 4$ and $F(1000) = 1069$.

\textbf{Find $F(10^{12})$.}

\section{Derivation}

\subsection{Substitution}

From $\tfrac{1}{x} + \tfrac{1}{y} = \tfrac{1}{n}$ with $x < y$ (hence $x > n$
and $y > n$), set $x = n + a$, $y = n + b$ with $0 < a < b$. Then:
\begin{align*}
\frac{1}{n+a} + \frac{1}{n+b} &= \frac{1}{n} \\
n(n+b) + n(n+a) &= (n+a)(n+b) \\
n^2 &= ab.
\end{align*}

So each solution corresponds to an ordered factorization $n^2 = ab$ with $a < b$
(equivalently $a < n$), subject to $y = n + b \leq L$.

\subsection{Coprime Parametrization}

\begin{proposition}
Every solution $(x, y, n)$ is uniquely represented as
\[
x = k \, u \, (u+v), \quad y = k \, v \, (u+v), \quad n = k \, u \, v,
\]
where $k \geq 1$, $1 \leq u < v$, and $\gcd(u, v) = 1$.
\end{proposition}

\begin{proof}
Given $n^2 = ab$ with $a < b$, write $n = k \cdot m$ where $m^2 \mid n^2$ extracts
the coprime structure. Specifically, let $d = \gcd(a, b)$, $a = d u^2 / \gcd(u,v)^2$...

More directly: set $g = \gcd(a, n)$. Write $a = g \alpha$, $n = g \gamma$ where $\gcd(\alpha, \gamma) = 1$. Then $ab = n^2$ gives $g \alpha \cdot b = g^2 \gamma^2$, so $b = g \gamma^2 / \alpha$. For $b$ to be an integer, $\alpha \mid g$. Set $g = \alpha k$, then $a = \alpha^2 k$, $n = \alpha \gamma k$, $b = \gamma^2 k$.

With $u = \alpha$, $v = \gamma$: $\gcd(u, v) = \gcd(\alpha, \gamma) = 1$, $a < b$ gives $u < v$, and
\[
x = n + a = k\alpha\gamma + k\alpha^2 = ku(\underbrace{u+v}), \quad
y = n + b = k\alpha\gamma + k\gamma^2 = kv(u+v).
\]
The map $(a, b, n) \mapsto (u, v, k)$ is bijective. \qed
\end{proof}

\subsection{Summation Formula}

The constraint $y \leq L$ becomes $k \cdot v \cdot (u+v) \leq L$, giving
$k \leq \lfloor L / (v(u+v)) \rfloor$. Summing over all valid triples:

\begin{theorem}
\[
F(L) = \sum_{\substack{v=2}}^{V} \;\; \sum_{\substack{u=1 \\ \gcd(u,v)=1}}^{v-1} \left\lfloor \frac{L}{v(u+v)} \right\rfloor,
\]
where $V = \lfloor\sqrt{L}\rfloor$ (since $v(v+1) \leq L$ is needed for $u=1$).
\end{theorem}

\section{Computational Approach}

\subsection{Direct Enumeration}

For $L = 10^{12}$, $V \approx 10^6$. The outer loop runs $10^6$ iterations.
For each $v$, the inner loop runs over $u < v$ with $\gcd(u,v) = 1$, totaling
$\varphi(v)$ terms. The overall complexity is:
\[
\sum_{v=2}^{V} \varphi(v) \approx \frac{3V^2}{\pi^2} \approx 3 \times 10^{11}.
\]

In optimized C++ (fast GCD, branch prediction, early termination when
$v(u+v) > L$), this completes within minutes.

\subsection{Early Termination}

For each $v$, as $u$ increases, $v(u+v)$ increases. Once $v(u+v) > L$, we can
break. Since $u + v > v$, the effective upper bound on $u$ is
$\lfloor L/v \rfloor - v$, which for large $v$ is much smaller than $v$.

\subsection{Verification Data}

\begin{center}
\begin{tabular}{rl}
\toprule
$L$ & $F(L)$ \\
\midrule
15 & 4 \\
100 & 60 \\
1{,}000 & 1{,}069 \\
10{,}000 & 15{,}547 \\
100{,}000 & 203{,}931 \\
1{,}000{,}000 & 2{,}524{,}207 \\
\bottomrule
\end{tabular}
\end{center}

\section{Answer}

\[
\boxed{5435004633092}
\]

\end{document}
